% Propietario conservado: /Users/ruben/Documents/New project/output/EXTREMA_MEDIA_RAZON_RECIPROCIDAD_ES_EN_20260918/ARTICULO/ES/source/sections/alpha.tex
\section{Dérivation de la constante de structure fine}
\label{x_reciprocidad:sec:alpha}

La constante de structure fine reçoit dans ce développement une détermination arithmétique qui réunit les trois coordonnées régionales et le registre intégral dodécaphasique. L'opération combine ses blocs avec orientation fixée et transporte les quotients de la normalisation positionnelle. Le paramètre résultant exprime la compatibilité de cette différence orientée avec la prolongation des cylindres. Sa valeur et les données d'incidence utilisées dans sa construction resteront liées pendant le développement exceptionnel.

L'ordre interne est important. Les coordonnées \(\pi\), \(e\) et \(\varphi\), déjà construites dans la section précédente, sont des coordonnées
précédentes du même état arithmétique avec mémoire de trajectoire, et peuvent agir légitimement comme
entrées d'un lecteur ultérieur. Elles ne sont pas des constantes conventionnelles
injectées pour sélectionner alpha. La génération corrélée des quatre coordonnées partage une origine et
une compatibilité coinductive, sans que cela impose d'évaluer ses quatre
composantes en un seul pas.

Les deux voies expriment des fonctions mathématiques différentes du même paramètre.
La première réalise une clôture entière par division euclidienne en base
\(1000\). La seconde étudie l’équilibre analytique des vacances au moyen
de la résolvante \(E_{21/22}\) et de ses prolongements gradués. Une voie
détermine des blocs et des quotients ; l’autre détermine une racine par une
équation et un domaine d’unicité. Leur identification est formulée ensuite
dans les cylindres communs à toute profondeur. L’exposé distingue
ces trois résultats : construction positionnelle, détermination analytique
et compatibilité des deux.

\subsection{Première voie : détermination positionnelle par clôture entière}
\subsubsection{Domaine et orientation de la clôture entière}
\label{x_reciprocidad:subsec:alpha-dominio}

Soient \(P_j,E_j,\Phi_j\in\{0,\ldots,999\}\) les blocs canoniques des
trois coordonnées régionales, et prolongeons le registre dodécaphasique par
\[
 K_j=K_{1+((j-1)\bmod12)}.
\]
L'orientation du canal est \((1,1,-1)\). La contribution du registre dodécaphasique se soustrait selon la clôture dodécaphasique déjà typée. Pour \(B=1000\), l'opération locale distingue le bilan signé
\[
 Z_j:=P_j+E_j-\Phi_j-K_j
\]
de la somme qui incorpore le quotient entrant. La normalisation est
\begin{equation}
 s_j=Z_j+c_{j+1},\qquad
 A_j=s_j-B\left\lfloor\frac{s_j}{B}\right\rfloor,
 \qquad c_j=\left\lfloor\frac{s_j}{B}\right\rfloor.
 \label{x_reciprocidad:eq:alpha-recurrencia}
\end{equation}
Ainsi, \(A_j\in\{0,\ldots,B-1\}\), même lorsque \(s_j\) est
négatif. La retenue qui entre depuis la droite appartient au domaine de
l'opération ; elle ne peut pas être omise car le tableau s'imprime de gauche
à droite.

\begin{definition}[Clôture de profondeur finie]
\label{x_reciprocidad:def:alpha-finita}
Pour une profondeur \(N\) et une frontière entière \(t\), l'application
\(\mathcal C_N\) reçoit
\((P_{1:N},E_{1:N},\Phi_{1:N},K_{1:N},c_{N+1}=t)\) et applique
\eqref{x_reciprocidad:eq:alpha-recurrencia} dans l'ordre \(N,N-1,\ldots,1\). Sa sortie
est la paire de mots
\[
 (A_{1:N},c_{1:N+1}).
\]
\end{definition}

La division euclidienne ne possède pas un choix de signe résiduel : le reste se
fixe en \(\{0,\ldots,999\}\) et le quotient reste déterminé. Par
exemple, \(-431=569-1000\) et \(-1200=800-2\cdot1000\). Les
retenues négatives de la clôture ne sont pas des anomalies ; elles sont l'information
entière nécessaire pour maintenir l'égalité orientée.

\subsubsection{Fenêtre dodécaphasique et conditions de retenue}
\label{x_reciprocidad:subsec:alpha-ventana}

Les douze premières triades obtenues sont
\begin{align*}
 P={}&(141,592,653,589,793,238,462,643,383,279,502,884),\\
 E={}&(718,281,828,459,045,235,360,287,471,352,662,497),\\
 \Phi={}&(618,033,988,749,894,848,204,586,834,365,638,117).
\end{align*}
Dans le canal de fermeture, les cinq régions qui déterminent la coordonnée \(P\) conservent leurs antécédents distincts. Partager cette lecture ne les identifie pas en tant que régions ou en tant qu'histoires.

Avec frontière \(c_{13}=0\), la récurrence produit le tableau suivant.
Chaque entrée est montrée, la retenue entrante, le total avant normalisation,
le bloc obtenu et la retenue sortante. Le tableau est un calcul entier
reproductible, non une liste de blocs sélectionnés par leur ressemblance avec
un objectif.

\begin{table}[htbp]
\centering
\small
\setlength{\tabcolsep}{3.5pt}
\begin{tabular}{r|rrrr|rr|rr}
\hline
\(j\)&\(P_j\)&\(E_j\)&\(\Phi_j\)&\(K_j\)&\(c_{j+1}\)&\(s_j\)&\(A_j\)&\(c_j\)\\
\hline
1&141&718&618&234&0&7&007&0\\
2&592&281&033&543&0&297&297&0\\
3&653&828&988&140&\(-1\)&352&352&0\\
4&589&459&749&729&\(-1\)&\(-431\)&569&\(-1\)\\
5&793&045&894&659&\(-2\)&\(-717\)&283&\(-1\)\\
6&238&235&848&824&\(-1\)&\(-1200\)&800&\(-2\)\\
7&462&360&204&621&0&\(-3\)&997&\(-1\)\\
8&643&287&586&058&\(-1\)&285&285&0\\
9&383&471&834&914&\(-1\)&\(-895\)&105&\(-1\)\\
10&279&352&365&794&0&\(-528\)&472&\(-1\)\\
11&502&662&638&146&0&380&380&0\\
12&884&497&117&601&0&663&663&0\\
\hline
\end{tabular}
\caption{Clôture initiale avec frontière droite nulle. La dépendance se propage de la dernière ligne à la première.}
\label{x_reciprocidad:tab:alpha-ventana}
\end{table}

Le mot des retenues est
\[
 (c_1,\ldots,c_{13})=(0,0,0,-1,-1,-2,-1,0,-1,-1,0,0,0),
\]
et la lecture de la fenêtre est le rationnel
\begin{equation}
 \alpha_{12}^{(0)}
 =\frac{7297352569283800997285105472380663}{10^{36}}
 =0.007297352569283800997285105472380663.
 \label{x_reciprocidad:eq:alpha-ventana-cero}
\end{equation}
L'exposant rappelle la frontière choisie. Cette fenêtre est exacte dans son
domaine fini ; on ne l'identifie pas d'emblée à la troncature de la
section infinie. La différence entre les deux lectures deviendra visible lors du
transport de la retenue de la queue lointaine.

\subsubsection{Télescopage et réversibilité dans une fibre de frontière}
\label{x_reciprocidad:subsec:alpha-telescopia}

La forme locale de l'égalité est
\begin{equation}
 P_j+E_j-\Phi_j-K_j-A_j=Bc_j-c_{j+1}.
 \label{x_reciprocidad:eq:alpha-identidad-local}
\end{equation}
Le terme de retenue est un lien entre des positions contiguës. Si on le supprime coordonnée par coordonnée, on obtient une égalité fausse. Seule une évaluation qui conserve ses extrémités permet de le télescoper.

\begin{proposition}[Identité entière de profondeur arbitraire]
\label{x_reciprocidad:prop:alpha-telescopia}
Pour \(\iota_N(z)=\sum_{j=1}^{N}z_jB^{N-j}\), toute exécution de
\(\mathcal C_N\) satisfait
\begin{equation}
 \iota_N(P)+\iota_N(E)-\iota_N(\Phi)-\iota_N(K)-\iota_N(A)
 =B^Nc_1-c_{N+1}.
 \label{x_reciprocidad:eq:alpha-telescopia}
\end{equation}
\end{proposition}

\begin{proof}
Multiplier \eqref{x_reciprocidad:eq:alpha-identidad-local} par \(B^{N-j}\) et ajouter  
produit du côté droit
\[
 \sum_{j=1}^{N}(Bc_j-c_{j+1})B^{N-j}=B^Nc_1-c_{N+1}.
\]
Chaque retenue intérieure apparaît deux fois avec un signe opposé ; seuls les deux extrêmes survivent.
\end{proof}

Pour le tableau ~\ref{x_reciprocidad:tab:alpha-ventana}, les deux extrémités sont nulles. Si
\(p_{12},e_{12},f_{12}\) sont les trois sommes fractionnaires tronquées,
\[
 \alpha_{12}^{(0)}=p_{12}+e_{12}-f_{12}-\kappa_{36}.
\]
La constante du registre dodécaphasique qui apparaît est \(\kappa_{36}\), non
\(\kappa_{\mathrm{per}}\). Le domaine de cette identité contient
des troncatures et une frontière finie ; sa forme n'autorise pas à effacer les deux
conditions en passant à la section complète.

L'inversion locale est aussi exacte :
\[
 K_j=P_j+E_j-\Phi_j+c_{j+1}-A_j-Bc_j.
\]
Étant donnés \(P,E,\Phi\), les blocs \(A\) et les retenues complètes,  
on restaure le registre dodécaphasique. La concaténation restaure en outre
\[
 \iota_N(K)=\iota_N(P)+\iota_N(E)-\iota_N(\Phi)-\iota_N(A)
             -B^Nc_1+c_{N+1}.
\]
Cette réversibilité est une propriété de la fibre avec données et frontières fixes. Elle ne doit pas devenir une nouvelle recette causale qui commence par alpha et fabrique ensuite le registre. La direction de construction reste encore registre signé, registre dodécaphasique, clôture et sortie.

\subsubsection{Prolongation et normalisation canonique de la section complète}
\label{x_reciprocidad:subsec:alpha-completa}

La prolongation conserve le registre dodécaphasique périodique et reçoit les suites
régionales complètes. Soient
\[
 p=\sum_{j\geq1}P_jB^{-j},\qquad
 e_0=\sum_{j\geq1}E_jB^{-j},\qquad
 f=\sum_{j\geq1}\Phi_jB^{-j}.
\]
Chaque série est une réalisation du caractère interne correspondant.  
Les parties entières déjà obtenues sont 3, 2 et 1, de sorte que  
\(p=\pi-3\), \(e_0=e-2\) et  
\(f=\varphi-1\).

\begin{proposition}[Existence et unicité de la normalisation complète]
\label{x_reciprocidad:prop:alpha-normalizacion}
La série
\begin{equation}
 y=p+e_0-f-\kappa_{\mathrm{per}}
   =\sum_{j\geq1}(P_j+E_j-\Phi_j-K_j)B^{-j}
 \label{x_reciprocidad:eq:alpha-serie}
\end{equation}
converge absolument. Dans le domaine des blocs obtenus,
\(0<y<10^{-2}\). Il existe une unique expansion canonique \((A_j)\) de
\(y\), sans queue finalement constante 999, et une unique suite entière
bornée \((c_j)\) qui satisfait \eqref{x_reciprocidad:eq:alpha-identidad-local} avec
\(c_1=0\). Ces retenues appartiennent à l'ensemble invariant
\(\mathcal C=\{-2,-1,0,1,2\}\).
\end{proposition}

\begin{proof}
Les coefficients de la série sont bornés par \(2(B-1)\) en module,  
de sorte qu'elle est dominée par une série géométrique. Si  
\(Z_j=P_j+E_j-\Phi_j-K_j\) et  
\(R_j(Z)=\sum_{r\geq0}Z_{j+r}B^{-r-1}\), chaque queue est la différence  
de deux sommes de deux queues canoniques et appartient à \((-2,2)\).  
Le premier coefficient est  
\(Z_1=141+718-618-234=7\) ; par conséquent  
\(y=(7+R_2(Z))/1000\in(0.005,0.009)\).

Soit \((A_j)\) l'expansion canonique de \(y\). Définissons
\[
 c_j=R_j(Z)-R_j(A).
\]
En multipliant l'égalité des deux séries par \(B^{j-1}\) et en séparant
ses premiers \(j-1\) termes on obtient
\[
 c_j=\sum_{r=1}^{j-1}(A_r-Z_r)B^{j-1-r}\in\mathbb Z,
\]
et \(c_1=0\). De plus, \(R_j(A)\in[0,1)\), donc
\(-3<c_j<2\), ce qui inclut toutes les retenues dans \(\mathcal C\).
Les identités de déplacement des queues donnent
\(Bc_j=Z_j-A_j+c_{j+1}\), qui est la récurrence requise.

Réciproquement, toute solution bornée à sortie canonique et
\(c_1=0\) se télescope à la limite en la même série \(y\). Le développement
canonique est unique et, une fois fixé, la formule des queues détermine
toutes les retenues. Enfin, si l'on applique une exécution finie avec
\(c_{j+1}\in\mathcal C\), on a \(-2000\leq s_j\leq2000\),
de sorte que son quotient euclidien appartient de nouveau à \(\mathcal C\).
\end{proof}

La proposition explicite la normalisation du lecteur ultérieur : elle ne remplace pas la génération des trois suites par une série conventionnelle introduite comme donnée. Une fois ces suites produites par l'état coinductif, leur combinaison et la frontière de la normalisation sont déterminées sans consulter alpha.

La procédure finie conservée dans le corpus propage les cinq
frontières de \(\mathcal C\) et conserve uniquement le préfixe commun à
toutes. Si deux profondeurs ont coalescé avant une frontière,
la récurrence déterministe de droite à gauche donne les mêmes blocs
à sa gauche. C'est la compatibilité avec la troncature. On ne
choisit pas une frontière parce qu'elle donne des chiffres favorables, et on ne déduit pas la
compatibilité à toute profondeur de l'observation d'une longue exécution.

La comparaison initiale montre pourquoi cette précaution importe. Dans
la prolongation stabilisée on obtient \(c_{13}^{\infty}=-1\), tandis
que la fenêtre isolée imposait \(c_{13}=0\). Ainsi,
\[
 884+497-117-601-1=662
\]
dans le dernier bloc du premier tour. Les onze premières triades ne
changent pas, mais la prolongation commence
\[
 (007,297,352,569,283,800,997,285,105,472,380,662,999,\ldots).
\]
Une apparition du bloc 999 ne viole pas la convention canonique : ce
qui est exclu est une queue qui soit finalement constante 999. La fenêtre
à frontière nulle se termine en 663 ; le troncage de la section complète
se termine en 662. Les deux affirmations décrivent des opérations différentes
et exactes. Elles ne doivent pas être fusionnées en une seule notation de projection.

\subsubsection{Définition opérationnelle et équation additive complète}
\label{x_reciprocidad:subsec:alpha-definicion}

\begin{definition}[Coordonnée alpha de la clôture]
\label{x_reciprocidad:def:alpha-operativa}
La voie entière détermine la coordonnée
\begin{equation}
 \alpha_A=\sum_{j\geq1}A_jB^{-j}
 =p+e_0-f-\kappa_{\mathrm{per}},
 \label{x_reciprocidad:eq:alpha-via-a}
\end{equation}
où \((A_j,c_j)\) est la normalisation canonique compatible de
\eqref{x_reciprocidad:eq:alpha-recurrencia}. Dans la réalisation scalaire des
coordonnées HMT,
\begin{equation}
 \boxed{\alpha_A=
 \pi+e-\varphi
 -4-\kappa_{\mathrm{per}}.}
 \label{x_reciprocidad:eq:alpha-identidad-completa}
\end{equation}
\end{definition}

L'entier quatre n'est pas calibré : il s'agit de
\(3+2-1\), combinaison des parties entières des trois coordonnées.
La base mille provient de la carte cubique et la période douze du registre dodécaphasique. Les
signes sont ceux du canal orienté et de la fermeture, et la frontière lointaine est
celle qu'impose la section compatible. Telles sont les origines effectives
des coefficients de l'équation.

À ce niveau, on peut décrire alpha comme la coordonnée scalaire du défaut de fermeture entre les trois réalisations régionales et le retour rationnel dodécaphasique, calculé avec une retenue compatible. Le terme \emph{défaut} désigne une différence algébrique déterminée, et non une perte arbitraire ni une imperfection du système. Cette description ne supprime pas la généalogie : sans les régions, le registre signé, l'orientation et la loi de retenue, l'égalité scalaire ne serait pas une exposition suffisante de la construction.

Il faut également maintenir séparés trois objets que la notation d'alpha peut accompagner dans des lecteurs ultérieurs : \(A\) est un mot de blocs ; le support négatif de pré-retenue s'obtient à partir du signe des totaux avant normalisation ; et un vecteur radial exceptionnel appartient à un codomaine d'incidence. Aucun d'eux n'est interchangeable avec la valeur réelle \(\alpha_A\).


% BEGIN OWNER /Users/ruben/Documents/New project/output/EXTREMA_MEDIA_RAZON_RECIPROCIDAD_ES_EN_20260918/ARTICULO/ES/source/sections/revision_alpha_notacion.tex
\paragraph{Bilan signé et somme avec quotient entrant.}
\label{x_reciprocidad:ampl-alfa-z-s-aclaracion}
La division positionnelle comporte deux étapes consécutives. Avec
\(K_j^{\mathrm{per}}=K_{1+((j-1)\bmod12)}\), qui désigne le
prolongement périodique également noté \(K_j\) dans
\eqref{x_reciprocidad:eq:alpha-recurrencia}, le bilan antérieur au transport du
quotient est
\[
 Z_j:=P_j+E_j-\Phi_j-K_j^{\mathrm{per}}.
\]
La variable $s_j$ de \eqref{x_reciprocidad:eq:alpha-recurrencia} désigne la somme qui
inclut déjà la contribution entrante, de sorte que
\[
 s_j=Z_j+c_{j+1},\qquad
 A_j=Z_j+c_{j+1}-1000c_j.
\]
La configuration incidencielle
$H^-_\alpha$ utilise les signes du bilan $Z_j$ de la fenêtre initiale ;
la détermination des chiffres utilise en outre les quotients de liaison
$c_{j+1}$ et $c_j$. Cette distinction conserve conjointement le support
signé et la normalisation du développement.

% END OWNER


% BEGIN OWNER /Users/ruben/Documents/New project/output/EXTREMA_MEDIA_RAZON_RECIPROCIDAD_ES_EN_20260918/ARTICULO/ES/source/sections/revision_alpha_frontera.tex
% Borrador focal de adición, NO demostración de la igualdad de dos vías.
% Puede incorporarse su proposición positiva sin alterar fórmulas actuales.
% La nota PROCEDENCIA_Y_ALCANCE.md delimita la compatibilidad no reconstruida.
\paragraph{Dépendance exacte à l'égard de la frontière de troncature.}
\label{x_reciprocidad:ampl-alfa-frontera-exacta}
Les coordonnées régionales et le registre dodécaphasique procèdent de la
génération commune APP--TRIT--TPK. Une fois produites, leur composition
positionnelle possède une dépendance explicite à l'égard de la profondeur
de lecture. Soient $B=1000$,
\[
 Z_j=P_j+E_j-\Phi_j-K_j^{\mathrm{per}},\qquad
 y_N=\sum_{j=1}^{N}Z_jB^{-j},\qquad
 R_{N+1}(Z)=\sum_{r\geq0}Z_{N+1+r}B^{-r-1}.
\]
La série complète de \eqref{x_reciprocidad:eq:alpha-serie} satisfait exactement
\[
 \alpha_A=y_N+B^{-N}R_{N+1}(Z),\qquad
 -2<R_{N+1}(Z)<2.
\]
Cette identité réunit la contribution de la fenêtre et celle de son
prolongement, toutes deux exprimées dans les mêmes coordonnées.

\begin{proposition}[Transport de la frontière et cylindre canonique]
\label{x_reciprocidad:ampl-alfa-prop-frontera}
Soit $A_j^{(N,t)}$ la sortie de la division positionnelle finie avec
$c_{N+1}=t\in\mathcal C=\{-2,-1,0,1,2\}$, et soit
$a_N^{(t)}=\sum_{j=1}^{N}A_j^{(N,t)}B^{-j}$. Pour les registres du
domaine considéré, on a
\begin{equation}
 \alpha_A-a_N^{(t)}
 =B^{-N}\bigl(R_{N+1}(Z)-t\bigr),
 \qquad
 |\alpha_A-a_N^{(t)}|<4B^{-N}.
 \label{x_reciprocidad:ampl-alfa-eq-frontera}
\end{equation}
Le quotient du prolongement complet est
$c_{N+1}^{\infty}=\lfloor R_{N+1}(Z)\rfloor$. Avec cette frontière,
\[
 0\leq\alpha_A-a_N^{(c_{N+1}^{\infty})}<B^{-N},
\]
de sorte que la fenêtre obtenue est le préfixe canonique d'ordre $N$.
\end{proposition}
\begin{proof}
Le télescopage de \eqref{x_reciprocidad:eq:alpha-identidad-local} donne
$a_N^{(t)}=y_N+tB^{-N}-c_1$. Comme $Z_1=7$, chaque état entrant
appartient à $\mathcal C$ et le premier dividende est compris entre $5$ et $9$,
on a $c_1=0$. La substitution de cette égalité dans la décomposition de la
série prouve \eqref{x_reciprocidad:ampl-alfa-eq-frontera}. La majoration résulte de
$R_{N+1}(Z)\in(-2,2)$ et de $t\in\mathcal C$.

Pour la normalisation canonique complète,
$c_{N+1}^{\infty}=R_{N+1}(Z)-R_{N+1}(A)$ et
$R_{N+1}(A)\in[0,1)$. Le caractère entier du quotient implique
$c_{N+1}^{\infty}=\lfloor R_{N+1}(Z)\rfloor$. La formule d'erreur
devient alors
$\alpha_A-a_N^{(c_{N+1}^{\infty})}=B^{-N}R_{N+1}(A)$, ce qui prouve
l'appartenance au cylindre canonique.
\end{proof}

La proposition identifie la signification de la frontière droite : c'est la
partie entière de la queue signée exprimée à l'échelle de la position
suivante. Le calcul sur les cinq frontières conserve cette information
sous la forme d'une famille finie de possibilités jusqu'à la stabilisation des blocs à
gauche. La contraction quantitative procède de la
profondeur et du transport de la queue ; elle a pour objet la normalisation des
coordonnées déjà générées.

\paragraph{Domaine gradué de la résolvante analytique.}
\label{x_reciprocidad:ampl-alfa-dominio-resolvente}
La composition analytique présentée à la
\S\ref{x_reciprocidad:subsec:alpha-resolvente} agit sur
$\operatorname{span}\{e_7,e_8,e_9\}$ et satisfait $Ne_9=0$. Son
évaluation produit exactement $T_{7:9}$ ; le degré $9$ est la dernière
composante de ce domaine. L'identité $N^3=0$ explique pourquoi la
résolvante s'exprime en trois termes. Pour décrire un
prolongement en degrés supérieurs, il faut spécifier, en outre, le
transport entre domaines gradués successifs et son action sur la
composante initiale de chaque domaine. La nilpotence locale et la
compatibilité globale remplissent ainsi des fonctions mathématiques distinctes.

% END OWNER



% BEGIN OWNER /Users/ruben/Documents/New project/output/EXTREMA_MEDIA_RAZON_RECIPROCIDAD_ES_EN_20260918/ARTICULO/ES/source/sections/alpha_dependencias.tex
\subsubsection{Détermination constructive et fonction du registre dodécaphasique}
\label{x_reciprocidad:alpha-ampl-dependencias}

La définition précédente permet de préciser ce que signifie déterminer alpha à partir
de la structure. Le domaine de l'opération contient quatre registres de
provenance connue : les trois suites régionales et le registre entier
de phase. La normalisation incorpore en outre une condition de frontière
compatible avec le prolongement. Une fois ces données fixées, la division
euclidienne détermine chaque reste et son quotient. La proposition
\ref{x_reciprocidad:prop:alpha-normalizacion} étend cette détermination au système
complet. L'indépendance à l'égard d'une valeur physique cible est donc
une propriété de l'ordre de construction : la valeur comparée à la
fin n'intervient dans la sélection d'aucun des arguments de cette
application.

Cette indépendance causale doit être distinguée de l'indépendance algébrique
entre nombres. La première s'établit en identifiant le domaine, les
opérations et la provenance des coefficients ; la seconde serait une
assertion relative à des relations polynomiales. Le résultat utilisé ici est le
premier. La construction détermine une coordonnée au moyen de registres
produits par APP--TRIT--TPK et permet de vérifier ultérieurement les
relations qu'elle satisfait. Cet ordre permet d'utiliser une coordonnée déjà
engendrée dans une opération ultérieure sans la transformer en paramètre externe.

Le registre $K$ possède une fonction plus précise que celle d'une correction
numérique ajoutée à une formule. Ses douze positions proviennent de la
transformation réversible du registre signé. L'équation
$U=\mathcal H_{12}K$ conserve la possibilité de revenir à ce registre,
tandis que les différences transversales et la somme totale permettent une autre
reconstruction. Dans l'évaluation périodique, l'ordre des positions
détermine le poids positionnel de chaque composante. Le même ensemble de
composantes, disposé dans un autre ordre, produit en général une autre coordonnée
rationnelle. L'origine de phase et l'orientation font donc partie de la
définition du nombre qui intervient dans \eqref{x_reciprocidad:eq:alpha-identidad-completa}.

La normalisation d'alpha compose cette information avec les régions. Son
identité locale distingue deux opérations successives. On forme d'abord le
bilan orienté des blocs ; ce bilan est ensuite représenté par
un reste canonique et un quotient reliant des positions contiguës. La première
opération conserve le signe du bilan et permet l'extraction ultérieure
d'un support. La seconde permet de former les préfixes de la coordonnée
réelle. Les deux lectures restent liées parce que l'état conserve les
données antérieures et postérieures à la normalisation. Le support signé et le
développement positionnel sont ainsi des réalisations différentes d'une composition
arithmétique explicite.

L'identité télescopique exprime le contenu global de cette composition.
Les quotients intérieurs s'annulent par paires lors de l'évaluation d'un préfixe,
tandis que leurs extrémités demeurent dans l'égalité. Cette annulation est une
propriété de la somme pondérée ; elle n'élimine pas les quotients de l'état et
ne permet pas de les supprimer avant d'effectuer la somme. La frontière droite d'une
fenêtre reçoit la contribution de la continuation. Pour cette raison,
l'évaluation d'une fenêtre isolée et la restriction d'une section
prolongée sont des opérations distinctes. Le passage du bloc final $663$ à
$662$, démontré précédemment, fournit une réalisation explicite de cette
différence sans altérer les blocs déjà stabilisés à sa gauche.

En ce sens, la bidirectionnalité possède une formulation arithmétique
concrète. La génération fournit des prolongements compatibles ; la
normalisation transmet l'information entière de leur frontière depuis les
positions de moindre poids vers celles de poids supérieur. L'équation conserve la
contribution des deux directions au moyen de ses indices et de ses conditions
aux limites. La réversibilité dans la fibre fixée se vérifie en reconstruisant
$K$ à partir des autres registres et des quotients complets. Cette
inversion certifie la cohérence de l'opération d'origine ; la direction
généalogique reste celle qui produit d'abord le registre puis
la coordonnée d'alpha.

La seconde détermination possède un contenu complémentaire. L'objet
immédiat y est un bilan de lacunes, et le paramètre apparaît comme
sa racine dans un domaine isolant. La voie positionnelle fournit une
normalisation canonique des registres ; la voie analytique étudie l'annulation
d'un opérateur. L'exposé des deux exige de conserver précisément
cette différence de fonctions. Leurs relations avec l'état commun doivent
s'exprimer au moyen des applications de raffinement et de la comparaison de
leurs réalisations complètes, comme le développe la sous-section
correspondante. L'accord d'une approximation finie constitue un
contrôle localisé ; l'identification des constructions complètes
possède le domaine et la portée de son propre énoncé.

Enfin, cette organisation explique la position d'alpha dans l'article.
Son équation apparaît après la génération régionale et la construction
de $K$, parce qu'elle utilise les deux résultats. La section électronique utilise
ensuite cette coordonnée comme un facteur de sa composition. Le prolongement
exceptionnel reçoit les registres et leurs supports, où l'information
d'incidence est conservée. Ces deux continuations reprennent une même
provenance et développent des fonctions mathématiques distinctes : évaluation d'une
section centrale et construction d'une configuration de frontière.
Le développement de chacune permet de suivre la structure à travers
l'équation, au lieu de limiter la comparaison à sa valeur scalaire.

% END OWNER


\subsection{Seconde voie : détermination analytique au moyen des vacances}
\subsubsection{Invariants du transport et résolvante graduée}
\label{x_reciprocidad:subsec:alpha-via-b}

La voie analytique reçoit des invariants de l'état orienté, et non la sortie
\(\alpha_A\). Sa troncature d'ordre six est
\begin{equation}
 P_6(x)=2\pi x-\frac74x^2
 +\frac{x^3}{2\pi}
 +\frac{x^4}{20}-\frac{2x^5}{21}-\frac{x^6}{46},
 \label{x_reciprocidad:eq:alpha-p6}
\end{equation}
et la condition de vacance utilise
\[
 \delta=\log_{10}\frac{10}{9}.
\]
Le rapport \(10/9\) compare la carte décimale aux neuf positions APP de la construction nonadique. Son logarithme appartient à cette lecture ultérieure de compactification ; ce n'est pas une valeur cible d'alpha.

La généalogie des coefficients doit être conservée complète, mais avec sa
force logique exacte. Les développements récents réunissent une signature marquée
de l'état. La matrice de synchronisation et le recensement associé sont
\[
 D_{\mathrm{syn}}=\begin{pmatrix}85&112\\41&52\end{pmatrix},
 \qquad N_9=43.
\]
De son déterminant on extrait
\[
 a=-\frac{\det D_{\mathrm{syn}}}{N_9}
   =\frac{172}{43}=4.
\]
L'horloge induite sur les trous courts a un spectre secondaire \(\{6,7\}\) ; son extrémité supérieure donne \(b=7\). La cardinalité tritique \(m=3\) détermine
\[
 t_*=mb=21,\qquad k_*=mb-1=20.
\]
La lecture de vacance dans la carte mille donne
\[
 v_-=\lfloor1000\delta\rfloor=45,\qquad
 v_+=\lceil1000\delta\rceil=46.
\]
Ces entiers sont des antécédents typés de l'évaluation de la troncature. Ils ne sont obtenus ni à partir de sa racine ni à partir d'une approximation métrologique de celle-ci.

La branche inférieure coïncide en outre avec le déterminant du bloc APP
\[
 B_c=\begin{pmatrix}7&2\\2&7\end{pmatrix},\qquad\det B_c=45.
\]
Le rang transversal du registre du chapitre précédent donne
\(r_{\mathrm{dod}}=11\), et la demi-couronne de phase fournit
\(f_{\mathrm{ph}}=8\binom62=120\). En particulier,
\(45r_{\mathrm{dod}}=495\).

Avec \(\omega=2\pi\), la signature
\[
 \mathcal I_9=(\omega,m,a,b,k_*,t_*,v_-,v_+,f_{\mathrm{ph}},r_{\mathrm{dod}})
\]
alimente le lecteur gradué
\begin{align}
 \mathfrak E_9(\mathcal I_9)(x)
 ={}&\omega x-\frac ba x^2+\omega^{-1}x^3
 +\frac{x^4}{k_*}-\frac{2x^5}{t_*}-\frac{x^6}{v_+}\notag\\
 &-\frac{x^7}{f_{\mathrm{ph}}}+\frac{x^8}{v_-}
 +\frac{2x^9}{v_-r_{\mathrm{dod}}}.
 \label{x_reciprocidad:eq:alpha-firma-coeficientes}
\end{align}
La substitution restitue les coefficients de \(P_6\) et de sa
prolongation d'ordre neuf. Le couple \(45/46\) traverse la séparation
entre les deux : la branche supérieure occupe le degré six et la branche inférieure réapparaît
aux degrés huit et neuf.

La factorisation de \eqref{x_reciprocidad:eq:alpha-firma-coeficientes} est formulée
pour un état qui conserve origine, ordre des canaux et orientation. La
naturalité par rapport aux morphismes qui préservent ces données n’équivaut pas à
l’unicité de tout lecteur après leur oubli. Le domaine de cette
factorisation est déterminé par la signature
\eqref{x_reciprocidad:eq:alpha-firma-coeficientes} ; le prolongement à tous les degrés
est construit dans la section~\ref{x_reciprocidad:subsec:alpha-publicaciones-completas},
avec sa classe de récurrence et sa dépendance à la précoordonnée de l’état
expressément définies.

\subsubsection{Résolvante nilpotente aux degrés 7, 8 et 9}
\label{x_reciprocidad:subsec:alpha-resolvente}

La première prolongation admet une construction opératoire finie
spécialement explicite. Sur l'espace à base \(e_7,e_8,e_9\),
soit
\[
 Ne_7=-\frac83e_8,\qquad Ne_8=\frac2{11}e_9,\qquad Ne_9=0,
 \qquad v_7=-\frac1{120}e_7.
\]
Le symbole \(N\) désigne ici un opérateur nilpotent, et non une profondeur de troncature. Son domaine et sa graduation sont fixés avant d'évaluer le polynôme en un nombre.

\begin{lemma}[Prolongement nilpotent exact]
\label{x_reciprocidad:lem:alpha-resolvente}
On vérifie \(N^3=0\) et
\[
 (I-N)^{-1}v_7
 =-\frac{e_7}{120}+\frac{e_8}{45}+\frac{2e_9}{495}.
\]
L'évaluation \(e_n\mapsto x^n\) produit
\begin{equation}
 T_{7:9}(x)=-\frac{x^7}{120}+\frac{x^8}{45}+\frac{2x^9}{495}.
 \label{x_reciprocidad:eq:alpha-t79}
\end{equation}
\end{lemma}

\begin{proof}
L'opérateur augmente le degré et s'annule sur le dernier vecteur de la base,
de sorte que son cube est nul. L'inverse est la résolvante finie
\(I+N+N^2\). De plus,
\[
 Nv_7=\frac1{45}e_8,\qquad N^2v_7=\frac2{495}e_9.
\]
Additionner les trois termes et évaluer donne la formule.
\end{proof}

Ici, les dénominateurs 120, 45 et 495 ne sont pas des ajustements décimaux : ils proviennent de la demi-couronne, du déterminant APP et de son produit par le rang dodécaphasique. Les transferts \(-8/3\) et \(2/11\), ainsi que le signe et le degré du germe, appartiennent au lecteur typé. Conserver seulement la symétrie du cycle et remplacer cet opérateur par un autre opérateur symétrique ne conserve pas nécessairement les coefficients.

\subsubsection{Racine de multiplicité un de la troncature d'ordre neuf}
\label{x_reciprocidad:subsec:alpha-jet}

Définissons
\[
 P_9=P_6+T_{7:9},\qquad E_{21/22}^{(9)}(x)=P_9(x)-\delta,
 \qquad I_\alpha=(0,10^{-2}).
\]
Cet opérateur fini a un résultat propre d'existence et d'unicité.

\begin{proposition}[Racine de multiplicité un de la troncature]
\label{x_reciprocidad:prop:alpha-raiz-jet}
L'opérateur \(E_{21/22}^{(9)}\) a une unique racine \(\xi_9\) dans \(I_\alpha\), et cette racine est simple. On obtient l'isolement rationnel
\begin{equation}
 \frac{729735256928380099}{10^{20}}
 <\xi_9<\frac{729735256928380100}{10^{20}}.
 \label{x_reciprocidad:eq:alpha-intervalo-jet}
\end{equation}
\end{proposition}

\begin{proof}
Les bornes des coordonnées internes
\(3.14<\pi<3.15\) et
\(0.045<\delta<0.046\) sont suffisantes pour l'existence. En zéro,
\(E_{21/22}^{(9)}(0)=-\delta<0\). En \(10^{-2}\), la contribution
linéaire moins les termes négatifs est supérieure à \(0.0626\), donc la
valeur de l'opérateur est positive.

Pour \(0\leq x\leq10^{-2}\), en éliminant les termes positifs de
la dérivée on obtient
\[
 (E_{21/22}^{(9)})'(x)
 \geq2\pi-\frac72\cdot10^{-2}
 -\frac{10}{21}\cdot10^{-8}-\frac6{46}\cdot10^{-10}
 -\frac7{120}\cdot10^{-12}>6.24.
\]
La continuité donne l'existence ; la borne de la dérivée donne l'unicité et la simplicité. L'intervalle plus étroit \eqref{x_reciprocidad:eq:alpha-intervalo-jet} résulte d'une évaluation rationnelle dirigée. Soient \(x_-\) et \(x_+\) ses extrémités. Pour calculer \(\delta=\log(10/9)/\log 10\), on utilise, pour \(z>1\),
\[
 L_N(z)=2\sum_{j=0}^{N-1}\frac{y^{2j+1}}{2j+1},\qquad
 y=\frac{z-1}{z+1},\qquad
 0<\log z-L_N(z)<\frac{2y^{2N+1}}{(2N+1)(1-y^2)}.
\]
Le reste est borné au moyen de la série géométrique des termes positifs.
On applique cette formule à \(z=10/9\) et \(z=3\), avec \(N=520\),
en utilisant \(\log 10=\log(10/9)+2\log3\). L’évaluation de
\(\pi\) utilise le cylindre régional engendré à mille chiffres, en incluant
sa largeur \(10^{-1000}\). L’arithmétique des intervalles dans
\eqref{x_reciprocidad:eq:alpha-firma-coeficientes} donne
\[
 -\frac{4559}{10^{23}}<E_{21/22}^{(9)}(x_-)< -\frac{4558}{10^{23}},
 \qquad
 \frac{1698}{10^{23}}<E_{21/22}^{(9)}(x_+)<\frac{1699}{10^{23}}.
\]
Ces signes et la monotonie prouvée établissent l’isolement. Les extrémités
sont des résultats de l’évaluation de l’opérateur ; ses coefficients proviennent de
la signature préalablement construite.
\end{proof}

Le résultat concerne \(\xi_9\). Il n'autorise pas à écrire \(\xi_9=\alpha_A\) comme une égalité exacte de nombres complets. La racine d'une troncature peut partager un préfixe avec la racine de l'opérateur prolongé et différer à une position ultérieure. La simplicité permet de contrôler le déplacement de la racine sous perturbations ; elle n'annule pas ces perturbations.

\subsection{Compatibilité des deux déterminations à toute profondeur}
\label{x_reciprocidad:subsec:alpha-limite-dos-vias}

\begin{remark}[Troncature, prolongement et compatibilité]
Les opérations des degrés \(6\) à \(9\) déterminent un jet fini. La normalisation positionnelle détermine la coordonnée complète au moyen de sa suite compatible de retenues. La proposition suivante exprime le critère général d'identification par cylindres. La section~\ref{x_reciprocidad:subsec:alpha-publicaciones-completas} construit ensuite une carte analytique corrélée de tous les degrés, avec récurrence, convergence uniforme, unicité de la racine et compatibilité projective. La construction reçoit la coordonnée partagée de l'état et la classe déclarée de prolongements ; elle conserve ainsi la différence entre une représentation analytique de cette coordonnée et une sélection indépendante.
\end{remark}

La voie analytique complète s’organise en une famille de lecteurs
gradués compatibles
\[
 E_{21/22}^{(6)}\longleftarrow E_{21/22}^{(9)}
 \longleftarrow E_{21/22}^{(12)}\longleftarrow\cdots.
\]
Le symbole
\[
 E_{21/22}=\varprojlim_m E_{21/22}^{(6+3m)}
\]
désigne ce lecteur complet, avec ses transports et ses troncatures, non une
série formelle choisie librement pour la faire passer par \(\alpha_A\).
La seconde voie détermine sa racine positive unique dans \(I_\alpha\),
notée \(\alpha_B\), avec l’unicité et la compatibilité du
lecteur complet démontrées dans la section~\ref{x_reciprocidad:subsec:alpha-publicaciones-completas}.

Les démonstrations de la section~\ref{x_reciprocidad:subsec:alpha-publicaciones-completas}
établissent que les deux systèmes déterminent, à chaque profondeur, le même
cylindre survivant \(C_N^\alpha\), et que leurs diamètres tendent vers
zéro. Il convient d’exposer séparément la conséquence mathématique de cette
compatibilité.

\begin{proposition}[Identification des limites compatibles]
\label{x_reciprocidad:prop:alpha-dos-limites}
Si les valeurs des voies entière et analytique appartiennent aux mêmes cylindres \(C_N^\alpha\) pour tout \(N\), et
\(\operatorname{diam}C_N^\alpha\to0\), alors
\[
 \alpha_A=\alpha_B=\alpha.
\]
\end{proposition}

\begin{proof}
Pour tout \(N\), la distance entre les deux valeurs est bornée
par \(\operatorname{diam}C_N^\alpha\). Faire tendre \(N\) vers l'infini
donne une distance nulle. L'existence des sections compatibles assure
qu'il ne s'agit pas de l'intersection vide d'un système d'approximations.
\end{proof}

Cette preuve identifie la conséquence des carrés de compatibilité.
Leur construction à une profondeur arbitraire figure dans la
section~\ref{x_reciprocidad:subsec:alpha-publicaciones-completas}. Les évaluations
finies confrontent des réalisations de cette construction ; la récurrence
et ses bornes uniformes établissent le passage à la limite.

La différence peut s'exprimer algébriquement. Soit \(F=\mathbb Q(\pi,\delta)\) et soit \(j_9:F[[x]]\to F[x]/(x^{10})\) la troncature. Alors
\[
 \ker j_9=x^{10}F[[x]].
\]
Tous les opérateurs
\[
 E_h(x)=E_{21/22}^{(9)}(x)+x^{10}h(x)
\]
ont la même troncature d'ordre neuf. Cependant, en prenant
\(h=0\) et \(h=1/729\),
\[
 E_{1/729}(\xi_9)=\frac{\xi_9^{10}}{729}>0,
\]
de sorte qu'ils ne partagent pas cette racine. L'observation ne permet pas de choisir une queue étrangère au transport HMT ; elle démontre précisément pourquoi la queue générée et sa compatibilité doivent être conservées. Une série formelle arbitraire ne suffit pas non plus pour parler de racines analytiques sans son domaine de convergence ou sans la réalisation du système de cylindres.

% BEGIN OWNER /Users/ruben/Documents/New project/output/EXTREMA_MEDIA_RAZON_RECIPROCIDAD_ES_EN_20260918/ARTICULO/ES/source/sections/alpha_publicacion_completa.tex
\subsection{Compatibilité entre les réalisations dodécaphasique et analytique d’alpha}
\label{x_reciprocidad:subsec:alpha-publicaciones-completas}

La compatibilité est établie à toute profondeur. La publication
dodécaphasique réversible et la représentation analytique de la
coordonnée d'alpha
\(E_{21/22}\) sont deux publications internes de codomaines distincts,
produites à partir du même état enrichi APP--TRIT--TPK. Le jet
\(E_{21/22}^{(9)}\) est un témoin fini exact, mais ne détermine pas à lui seul
la queue. La complétion construite ensuite ne se déduit pas du jet
seul : elle reçoit la coordonnée partagée de l'état et publie son image
analytique par une récurrence explicite. Elle ne constitue donc pas une
sélection causale indépendante ; elle constitue une représentation analytique
typée du même point généré.

Cette séparation conserve l'ordre causal
\[
 \mathrm{APP}\longrightarrow\mathrm{TRIT}\longrightarrow\mathrm{TPK}
 \longrightarrow\mathcal X_{\alpha}^{\mathrm{enr}}
 \longrightarrow (P,E,\Phi,U_{12}^{\mathrm{sgn}},K)
 \longrightarrow\alpha_{\mathrm{HMT}}.
\]
Les coordonnées \(P,E,\Phi\) et le registre \(K\) sont des sorties antérieures du
même état enrichi. Leur utilisation dans la clôture d'alpha est donc une
composition interne ultérieure et non une injection de constantes
conventionnelles.

Pour \(r\in6\mathbb N\), \(r\ge36\), soit
\(\widehat{\mathcal X}^{\mathrm{enr},\alpha}_r\) le secteur des états
tronqués qui prolongent le germe interne d'alpha fixé dans \(R_{36}\). On
écrit
\[
 x_r^\alpha=(\widetilde x_r,\chi_\alpha,
              \varepsilon_r^\alpha,N_r^\alpha),
\]
où \(\chi_\alpha\) est le caractère interne du secteur, et non la valeur scalaire
publiée. Chaque état conserve résidu, quotient, retenue, feuillet, orientation,
phase, route, cylindre, frontière, signature, registre, mémoire, survie,
provenance et profondeur. Si
\[
 d_r^\alpha=\operatorname{Dig}_{729}(\varepsilon_r^\alpha),\qquad
 \varepsilon_{r+6}^\alpha=729\varepsilon_r^\alpha-d_r^\alpha,
 \qquad N_{r+6}^\alpha=729N_r^\alpha+d_r^\alpha,
\]
la transition effective est
\[
 \mathcal U_r^\alpha=
 \operatorname{Upd}_r^\alpha\circ
 \operatorname{Tra}_r^\alpha\circ
 \operatorname{Sel}_r^\alpha,
 \qquad
 \operatorname{Ext}^{\alpha,\to}_{r+6,r}
 =\operatorname{graph}(\mathcal U_r^\alpha).
\]
La fonctionnalité de ces extensions, une fois le germe \(R_{36}\) fixé,
produit une unique section compatible : la phase revient tandis que la mémoire et
la profondeur avancent. Le domaine global des deux déterminations est
\begin{equation}
 \mathcal X_\alpha^{\mathrm{enr}}:=\widehat\Omega_\alpha
 :=\varprojlim_{r\ge36}
 \bigl(\widehat{\mathcal X}^{\mathrm{enr},\alpha}_r,
       \operatorname{Ext}^{\alpha,\to}_{r+6,r}\bigr).
 \label{x_reciprocidad:eq:alpha-dominio-enriquecido-completo}
\end{equation}

Dans cette section, \(\kappa_{\mathrm{per}}\) est la lecture périodique \(\kappa_{\mathrm{ag}}\) du registre déjà construit ; les deux notations désignent la même fraction dans la carte de base mille fixée.

\subsubsection{La clôture dodécaphasique comme section complète}

Pour \(N\geq1\), posons
\[
 \mathcal W_N:=\{0,\ldots,999\}^{N},
 \qquad
 \mathsf T_{\alpha,N}(\mathsf s):=(A_1,\ldots,A_N),
 \qquad \mathsf s\in\mathcal X_{\alpha}^{\mathrm{enr}},
\]
où les blocs \(A_j\) s'obtiennent au moyen de
\eqref{x_reciprocidad:eq:alpha-recurrencia} avec la frontière compatible
\(c_{N+1}^{\infty}=\lfloor R_{N+1}(Z)\rfloor\). Si \(N'\geq N\), soit
\(\rho_{N',N}:\mathcal W_{N'}\to\mathcal W_N\) la restriction aux
\(N\) premiers blocs.

\begin{theorem}[Sortie dodécaphasique complète et compatibilité projective]
\label{x_reciprocidad:thm:alpha-salida-dodecafasica-completa}
La famille \((\mathsf T_{\alpha,N})_{N\geq1}\) est compatible :
\begin{equation}
 \rho_{N',N}\circ\mathsf T_{\alpha,N'}
 =\mathsf T_{\alpha,N}
 \qquad(N'\geq N).
 \label{x_reciprocidad:eq:alpha-via-a-proyectiva-exacta}
\end{equation}
Ses cylindres canoniques
\begin{equation}
 C_N^A(\mathsf s):=
 \left[
  \sum_{j=1}^{N}A_j1000^{-j},
  \sum_{j=1}^{N}A_j1000^{-j}+1000^{-N}
 \right)\cap I_\alpha
 \label{x_reciprocidad:eq:alpha-cilindros-via-a-exactos}
\end{equation}
sont emboîtés, ont un diamètre au plus égal à \(1000^{-N}\) et satisfont
\begin{equation}
 \bigcap_{N\geq1}C_N^A(\mathsf s)
 =\{\alpha_A(\mathsf s)\},
 \qquad
 \alpha_A
 =\pi_{\mathrm{HMT}}+e_{\mathrm{HMT}}-\varphi_{\mathrm{HMT}}
  -4-\kappa_{\mathrm{per}}.
 \label{x_reciprocidad:eq:alpha-salida-a-completa}
\end{equation}
En conséquence, la première voie publie une unique section infinie, que
nous notons \(\alpha_{\mathrm{HMT}}:=\alpha_A\).
\end{theorem}

\begin{proof}
La proposition~\ref{x_reciprocidad:prop:alpha-normalizacion} fournit l'unique développement canonique et l'unique suite bornée de retenues. L'identité locale \eqref{x_reciprocidad:eq:alpha-identidad-local} et son télescopage \eqref{x_reciprocidad:eq:alpha-telescopia} conservent la frontière distante ; la proposition \ref{x_reciprocidad:ampl-alfa-prop-frontera} prouve que restreindre une exécution prolongée redonne exactement le préfixe précédent. Cela donne \eqref{x_reciprocidad:eq:alpha-via-a-proyectiva-exacta}. L'appartenance de la somme complète à chaque cylindre découle de l'identité des queues, et \(\operatorname{diam}C_N^A\leq1000^{-N}\to0\) ; l'intersection contient un seul point. Enfin, \eqref{x_reciprocidad:eq:alpha-identidad-completa} identifie ce point à l'expression de \eqref{x_reciprocidad:eq:alpha-salida-a-completa}. Tous ses coefficients et signes sont fixés avant la publication d'alpha.
\end{proof}

La fenêtre \(\alpha_{12}^{(0)}\), fermée avec \(c_{13}=0\), reste un
calcul fini exact ; elle ne doit pas être confondue avec la restriction de la section
complète, dont la frontière compatible satisfait \(c_{13}^{\infty}=-1\). En
particulier,
\begin{equation}
 \rho_{36,12}(\alpha_{A,36})=\alpha_{A,12}
 \label{x_reciprocidad:eq:alpha-rho-36-12-proyectiva}
\end{equation}
est une égalité de mots compatibles, et non une égalité entre des mots de
longueurs différentes ni une identification de la fenêtre isolée à la
limite.



% BEGIN OWNER /Users/ruben/Documents/New project/output/EXTREMA_MEDIA_RAZON_RECIPROCIDAD_ES_EN_20260918/ARTICULO/ES/source/sections/alpha_carta_analitica_completa.tex
% Representación analítica correlacionada de alfa: acción, convergencia e
% intervalos.
\subsubsection{Représentation analytique de la coordonnée d'alpha}
\label{x_reciprocidad:subsec:alpha-lector-completo-rev08}

Le prolongement analytique se spécifie au moyen de tous ses coefficients.
Soit \(\mathsf s\in\mathcal X_\alpha^{\mathrm{enr}}\). Avant
la normalisation décimale, le même état publie déjà la précoordonnée
\begin{equation}
 a(\mathsf s):=p(\mathsf s)+e_0(\mathsf s)-f(\mathsf s)
                 -\kappa_{\mathrm{per}}(\mathsf s)\in I_\alpha,
 \label{x_reciprocidad:eq:alpha-precoordenada-comun}
\end{equation}
avec \(p=\pi_{\mathrm{HMT}}-3\),
\(e_0=e_{\mathrm{HMT}}-2\) et
\(f=\varphi_{\mathrm{HMT}}-1\). De façon équivalente,
\(a=\pi_{\mathrm{HMT}}+e_{\mathrm{HMT}}-
\varphi_{\mathrm{HMT}}-4-\kappa_{\mathrm{per}}\). On n'introduit pas ici
\(\alpha\) conventionnel : \(p,e_0,f\) sont les trois caractères régionaux déjà
générés et \(\kappa_{\mathrm{per}}\) est la lecture rationnelle du registre
terminal. La clôture entière applique à
\eqref{x_reciprocidad:eq:alpha-precoordenada-comun} la normalisation de
\eqref{x_reciprocidad:eq:alpha-recurrencia} ; la carte qui suit publie sur la même
précoordonnée une représentation analytique compatible. Son graphe de
dépendance reçoit \((p,e_0,f,\kappa_{\mathrm{per}})\) directement : il ne
contient pas d'arête \(\alpha_A\to\lambda\), bien que l'égalité démontrée plus
bas permette ensuite d'identifier les deux noms de publication.

Écrivons \(q=1/729\), soit \(F_{\mathsf s}\) le corps ordonné engendré par
les coordonnées de \(\mathsf s\), et posons
\(E_{9,\mathsf s}(X)=E_{21/22}^{(9)}(\mathsf s;X)\). Nous définissons le coefficient
de liaison
\begin{equation}
 \lambda(\mathsf s):=
 -\frac{E_{9,\mathsf s}(a(\mathsf s))
         \bigl(1-q\,a(\mathsf s)^3\bigr)}{a(\mathsf s)^{10}}.
 \label{x_reciprocidad:eq:alpha-lambda-estado}
\end{equation}
Tous ses arguments appartiennent à l'état avant la reconnaissance métrologique. En particulier, \(\lambda\) ne s'obtient pas en ajustant une chaîne décimale cible. L'équation~\eqref{x_reciprocidad:eq:alpha-lambda-estado} doit être lue avec son type exact : une fois que la première publication de l'état a produit \(a(\mathsf s)\), elle fixe l'unique prolongement de la classe géométrique définie ci-après. Elle ne constitue pas une deuxième sélection indépendante de \(a(\mathsf s)\), mais une représentation analytique corrélée du même état enrichi.

En effet, pour \(\mu\in F_{\mathsf s}\), posons
\[
 E_{\mathsf s,\mu}(X)
 :=E_{9,\mathsf s}(X)+\mu\frac{X^{10}}{1-qX^3}.
\]
Comme \(0<a(\mathsf s)<10^{-2}\) et \(qa(\mathsf s)^3<1\), on a
\begin{equation}
 E_{\mathsf s,\mu}(a(\mathsf s))=0
 \quad\Longleftrightarrow\quad
 \mu=\lambda(\mathsf s).
 \label{x_reciprocidad:eq:alpha-unicidad-lambda-clase-geometrica}
\end{equation}
Ainsi, le germe de coefficients ne contient pas de liberté une fois fixés l'état
et la classe de prolongements \(q\)-géométriques. Cette unicité est relative à la
classe déclarée ; elle n'affirme pas que le jet d'ordre neuf sélectionne par lui-même une
queue parmi toutes les séries analytiques possibles.

L'action locale sur des blocs de trois coefficients est
\begin{equation}
 \mathcal C_\alpha:F_{\mathsf s}^{3}\longrightarrow F_{\mathsf s}^{3},
 \qquad \mathcal C_\alpha(u,v,w)=q(u,v,w),
 \qquad c^{(0)}=(\lambda(\mathsf s),0,0).
 \label{x_reciprocidad:eq:alpha-accion-coeficientes}
\end{equation}
Par conséquent, pour tout \(n\geq0\),
\begin{equation}
 b_{10+3n}=q^n\lambda(\mathsf s),\qquad
 b_{11+3n}=b_{12+3n}=0.
 \label{x_reciprocidad:eq:alpha-recurrencia-coeficientes}
\end{equation}
C'est la règle qui calcule chaque coefficient au-delà du degré neuf. Il ne
reste aucun paramètre libre à chaque prolongation.

Pour \(M\geq0\), soit
\begin{equation}
 E_{\mathsf s,M}(X):=E_{9,\mathsf s}(X)
 +\lambda(\mathsf s)\sum_{n=0}^{M}q^nX^{10+3n}.
 \label{x_reciprocidad:eq:alpha-truncamientos-completos}
\end{equation}
La fonction complète est
\begin{equation}
 E_{\mathsf s,\infty}(X)
 :=E_{9,\mathsf s}(X)
 +\lambda(\mathsf s)\frac{X^{10}}{1-X^3/729}.
 \label{x_reciprocidad:eq:alpha-funcion-completa-explicita}
\end{equation}

Pour formuler le système projectif sans types implicites, nous fixons \(d_M=12+3M\) et définissons la fibration de jets
\begin{equation}
 \mathfrak J_M:=
 \coprod_{\mathsf s\in\mathcal X_\alpha^{\mathrm{enr}}}
 \{\mathsf s\}\times F_{\mathsf s}[X]/(X^{d_M+1}).
 \label{x_reciprocidad:eq:alpha-espacio-jets-rev08}
\end{equation}
Si \(M'\ge M\), la restriction est
\begin{equation}
 \tau_{M',M}(\mathsf s,[P]_{d_{M'}})
 :=(\mathsf s,[P]_{d_M}).
 \label{x_reciprocidad:eq:alpha-truncamiento-jets-rev08}
\end{equation}

\begin{proposition}[Naturalité de la récurrence analytique]
\label{x_reciprocidad:prop:alpha-naturalidad-recurrencia-explicita}
Soit
\[
 \Gamma_M(\mathsf s)
 =\bigl(\mathsf s,[E_{\mathsf s,M}]_{d_M}\bigr).
\]
Si \(M'\geq M\), la troncature \(\tau_{M',M}\) satisfait
\begin{equation}
 \tau_{M',M}\circ\Gamma_{M'}=\Gamma_M,
 \qquad
 \Gamma_{E21}(\mathsf s)=\varprojlim_M\Gamma_M(\mathsf s).
 \label{x_reciprocidad:eq:alpha-naturalidad-explicita}
\end{equation}
La projection sur les trois nouveaux degrés est précisément
\(\mathcal C_\alpha^{M+1}(c^{(0)})\).
\end{proposition}

\begin{proof}
L'équation~\eqref{x_reciprocidad:eq:alpha-recurrencia-coeficientes} fixe le bloc de degrés \(10+3n,11+3n,12+3n\). Lors du passage de \(M\) à \(M+1\), seul le bloc suivant est ajouté ; le tronquer redonne littéralement le polynôme précédent. La compatibilité pour \(M'\geq M\) s'obtient par composition. Ainsi, la limite inverse n'est pas une section choisie du jet nu, mais l'orbite unique de \(c^{(0)}\) sous \(\mathcal C_\alpha\).
\end{proof}

\subsubsection{Convergence uniforme, racine simple et intervalles rationnels}
\label{x_reciprocidad:subsec:alpha-cotas-completas-rev08}

Les cylindres régionaux et le registre périodique fournissent des bornes au moyen
de l'arithmétique des intervalles rationnels. Pour chaque coordonnée
\[
 c\in\{\pi_{\mathrm{HMT}},e_{\mathrm{HMT}},\varphi_{\mathrm{HMT}}\},
\]
soit \(P_c\) son préfixe nonadique certifié de mille chiffres. La donnée utilisée
n'est pas le rationnel \(P_c\) comme s'il s'agissait de la valeur complète, mais le cylindre
semi-ouvert et sa fermeture extérieure
\[
 \mathcal C_c^{(1000)}=[P_c,P_c+\varepsilon),
 \qquad
 \overline{\mathcal C}_c^{(1000)}=[P_c,P_c+\varepsilon],
 \qquad \varepsilon=10^{-1000}.
\]
L'arithmétique dirigée opère de manière conservative avec ces fermetures. Si
\(P_\pi,P_e,P_\varphi\) désignent les trois préfixes, on obtient
\begin{equation}
\begin{aligned}
 a_-&:=P_\pi+P_e-(P_\varphi+\varepsilon)-4-
       \kappa_{\mathrm{per}},\\
 a_+&:=(P_\pi+\varepsilon)+(P_e+\varepsilon)-P_\varphi-4-
       \kappa_{\mathrm{per}},\\
 A_{1000}&:=[a_-,a_+],
 \qquad a(\mathsf s)\in[a_-,a_+]=A_{1000},
 \qquad \operatorname{diam}A_{1000}=3\varepsilon.
\end{aligned}
\label{x_reciprocidad:eq:alpha-propagacion-colas-rev08}
\end{equation}
L'extension naturelle par intervalles de
\eqref{x_reciprocidad:eq:alpha-lambda-estado} transporte ensuite \(A_{1000}\), le cylindre
de \(\pi_{\mathrm{HMT}}\) et la borne rationnelle dirigée de \(\delta\) vers un
intervalle fermé \(\Lambda_{1000}\ni\lambda(\mathsf s)\) ; aucune
des trois queues n'est figée. En particulier,
\begin{equation}
\begin{aligned}
 \frac{7297352569283800997285105472380662}{10^{36}}
 &<a(\mathsf s)<
 \frac{7297352569283800997285105472380663}{10^{36}},\\
 -8.711\mathbin{\times}10^{-6}
 &<\lambda(\mathsf s)<-8.709\mathbin{\times}10^{-6}.
\end{aligned}
 \label{x_reciprocidad:eq:alpha-cotas-precoordenada-lambda}
\end{equation}
Ces inégalités utilisent des préfixes certifiés avec leur borne rationnelle de
queue ; elles ne reçoivent pas de table d'alpha. Le vérificateur du paquet les reproduit
à partir des trois sorties nonadiques et de la carte terminale.

Soit
\[
 r:=\frac{10^{-6}}{729}<1.372\mathbin{\times}10^{-9}.
\]
Pour \(0\leq x\leq10^{-2}\), la queue postérieure à \(M\) satisfait
\begin{equation}
 \left|E_{\mathsf s,\infty}(x)-E_{\mathsf s,M}(x)\right|
 \leq
 8.711\mathbin{\times}10^{-26}\,\frac{r^{M+1}}{1-r}.
 \label{x_reciprocidad:eq:alpha-cota-cola-uniforme}
\end{equation}
De plus,
\begin{equation}
 \left|\frac{\mathrm d}{\mathrm dx}
  \left(\lambda\frac{x^{10}}{1-x^3/729}\right)\right|
 \leq 8.711\mathbin{\times}10^{-6}
 \left(\frac{10\,10^{-18}}{1-r}
 +\frac{3\,10^{-24}/729}{(1-r)^2}\right)
 <8.712\mathbin{\times}10^{-23}.
 \label{x_reciprocidad:eq:alpha-cota-derivada-cola}
\end{equation}
La queue des dérivées postérieure au bloc \(M\) admet en outre l'estimation dépendant de la profondeur
\begin{equation}
 \sup_{0\le x\le10^{-2}}
 \left|E'_{\mathsf s,\infty}(x)-E'_{\mathsf s,M}(x)\right|
 \le 8.711\mathbin{\times}10^{-24}\,r^{M+1}
 \left(\frac{13+3M}{1-r}+\frac{3r}{(1-r)^2}\right)
 \xrightarrow[M\to\infty]{}0.
 \label{x_reciprocidad:eq:alpha-cota-resto-derivadas-rev08}
\end{equation}

\begin{theorem}[Fonction limite et unique racine simple de la carte corrélée]
\label{x_reciprocidad:thm:alpha-raiz-completa-explicita}
La suite \(E_{\mathsf s,M}\) converge uniformément, ainsi que ses
dérivées, vers \(E_{\mathsf s,\infty}\) sur
\(\overline I_\alpha=[0,10^{-2}]\). On a
\begin{equation}
 E_{\mathsf s,\infty}(a(\mathsf s))=0,
 \qquad
 \inf_{x\in\overline I_\alpha}
 E'_{\mathsf s,\infty}(x)>6.239.
 \label{x_reciprocidad:eq:alpha-raiz-y-derivada-completas}
\end{equation}
Par conséquent, \(a(\mathsf s)\) est l'unique racine simple de la fonction complète dans \(I_\alpha\).
\end{theorem}

\begin{proof}
Le rapport de deux termes consécutifs de la queue est borné par \(r<1\),
ce qui donne \eqref{x_reciprocidad:eq:alpha-cota-cola-uniforme} ; la dérivation terme à
terme et la somme des deux séries géométriques donnent
\eqref{x_reciprocidad:eq:alpha-cota-derivada-cola} et
\eqref{x_reciprocidad:eq:alpha-cota-resto-derivadas-rev08}. La preuve du jet établit
\(E'_{9,\mathsf s}>6.24\) dans \(\overline I_\alpha\) ; par conséquent, la
perturbation complète maintient la dérivée au-dessus de \(6.239\).
Enfin, en substituant \eqref{x_reciprocidad:eq:alpha-lambda-estado} dans
\eqref{x_reciprocidad:eq:alpha-funcion-completa-explicita}, les deux termes évalués en
\(a(\mathsf s)\) s'annulent exactement. La fonction est strictement
croissante ; son unique racine est donc simple et coïncide avec cette précoordonnée.
\end{proof}

Pour rendre effective l'identification à toute profondeur, les évaluations
dirigées des extrémités certifient d'abord
\[
 E_{\mathsf s,\infty}(0)<0<
 E_{\mathsf s,\infty}(10^{-2}),
 \qquad
 E'_{\mathsf s,\infty}\geq\mu:=6239/1000>c:=6
 \quad\hbox{dans }[0,10^{-2}].
\]
Par continuité, il existe une racine dans l'intervalle et, par la deuxième inégalité, elle est unique. Ce n'est qu'après avoir établi cette existence, la monotonie et l'invariant selon lequel l'encadrement contient la racine que l'on applique le rattrapage par le théorème des accroissements finis. Nous partons de \(D_0=[0,10^{-2}]\). Si \(D_j=[r_j,s_j]\), soit \(m_j=(r_j+s_j)/2\), et soit \([F_j^-,F_j^+]\) une inclusion rationnelle dirigée de \(E_{\mathsf s,\infty}(m_j)\), obtenue en propageant les cylindres de \(\pi_{\mathrm{HMT}},e_{\mathrm{HMT}},\varphi_{\mathrm{HMT}}\), \(\delta\) et \(\lambda\), avec
\begin{equation}
 0\leq F_j^+-F_j^-\leq \frac{c}{4}\,(s_j-r_j).
 \label{x_reciprocidad:eq:alpha-anchura-evaluacion-rescate-rev08}
\end{equation}
Lorsque \(F_j^+<0\), on conserve la moitié droite ;
lorsque \(F_j^->0\), la moitié gauche. Si
\(F_j^-\leq0\leq F_j^+\), on ne cherche pas à décider si le milieu est
exactement la racine. Soit \(w_j=s_j-r_j\). La borne uniforme
\(E'_{\mathsf s,\infty}\geq c=6\) et le théorème des accroissements finis impliquent que
toute racine contenue dans \(D_j\) appartient à
\begin{equation}
 D_{j+1}=D_j\cap
 \left[m_j-\frac{w_j}{4},
             m_j+\frac{w_j}{4}\right].
 \label{x_reciprocidad:eq:alpha-rescate-derivada-rev08}
\end{equation}
En effet, soit \(\xi_j\) la racine et soit
\(y_j=E_{\mathsf s,\infty}(m_j)\in[F_j^-,F_j^+]\). Comme l'encadrement contient
zéro et satisfait \eqref{x_reciprocidad:eq:alpha-anchura-evaluacion-rescate-rev08}, on a
\(|y_j|\leq cw_j/4\). Le théorème des accroissements finis fournit
\(|m_j-\xi_j|\leq |y_j|/c\leq w_j/4\), ce qui prouve
\eqref{x_reciprocidad:eq:alpha-rescate-derivada-rev08}. L'argument inclut le cas
\(y_j=0\) : la racine est alors au milieu, mais l'algorithme n'a pas
besoin de reconnaître cette égalité pour la conserver.

Si l'encadrement disponible ne satisfait pas \eqref{x_reciprocidad:eq:alpha-anchura-evaluacion-rescate-rev08}, on raffine d'abord les cylindres sources et on répète l'évaluation ; on ne déclare aucun signe et on n'accepte pas de contraction insuffisante. Dans chacune des trois branches, le nouvel intervalle conserve la racine et a un diamètre au plus égal à \(w_j/2\). Pour chaque profondeur cible finie, le raffinement dirigé permet de satisfaire la borne de largeur sans comparer de manière non certifiée un nombre réel à zéro. La monotonie démontre par récurrence
\begin{equation}
 a(\mathsf s)\in D_{j+1}\subseteq D_j,\qquad
 \operatorname{diam}D_{j+1}\leq
 \tfrac12\operatorname{diam}D_j,
 \qquad \operatorname{diam}D_j\longrightarrow0,
 \label{x_reciprocidad:eq:alpha-biseccion-racional}
\end{equation}
où l'inégalité est une borne de contraction et non une égalité imposée.
Le certificat exécutable applique
cette règle à vingt-quatre niveaux en base mille et vérifie que la fermeture
extérieure complète \(A_{1000}\), et non sa seule extrémité inférieure, est contenue dans
chaque cylindre publié. Ses tests négatifs rejettent une enceinte trop
large, une enceinte désordonnée, un point autre que le milieu, un domaine sans
changement de signe et une fourchette dégénérée. La prolongation coinductive permet
de remplacer mille par toute profondeur requise.

Concrètement, pour \(S_N=1000^N\), le vérificateur calcule
\begin{equation}
 j_N^-:=\lfloor S_Na_-\rfloor,
 \qquad j_N^+:=\lfloor S_Na_+\rfloor
 \label{x_reciprocidad:eq:alpha-indices-cilindro-dirigidos-rev08}
\end{equation}
et ne publie le niveau que si \(j_N^-=j_N^+\). Cette égalité, qui inclut
l'extrémité supérieure de la fermeture extérieure, prouve
\(A_{1000}\subset[j_N^-/S_N,(j_N^-+1)/S_N)\) et évite d'attribuer à la cellule
gauche une extrémité située exactement sur sa frontière droite.

Soit \(C_N^A(\mathsf s)\) le cylindre entier de \eqref{x_reciprocidad:eq:alpha-cilindros-via-a-exactos}. Nous choisissons récursivement \(m(N)>m(N-1)\) de sorte que \(\operatorname{diam}D_{m(N)}<1000^{-N}/4\), et définissons
\begin{equation}
I_{B,N}:=D_{m(N)}\cap C_N^A(\mathsf s).
 \label{x_reciprocidad:eq:alpha-intervalos-b-completos}
\end{equation}
Nous définissons
\(\ell_N:=\inf I_{B,N}\) et \(u_N:=\sup I_{B,N}\). Tous deux sont rationnels ;
l'intervalle est non vide parce que les deux facteurs contiennent
\(a(\mathsf s)\). Le cylindre étant semi-ouvert, \(u_N\) peut ne pas
appartenir à \(I_{B,N}\). L'évaluation en cette extrémité s'entend au moyen
du prolongement continu de \(E_{\mathsf s,\infty}\) à la fermeture extérieure. Ainsi,
\begin{equation}
 I_{B,N+1}\subseteq I_{B,N}\subseteq C_N^A(\mathsf s),
 \qquad \operatorname{diam}I_{B,N}\longrightarrow0,
 \qquad E_{\mathsf s,\infty}(\ell_N)\leq0\leq
 E_{\mathsf s,\infty}(u_N).
 \label{x_reciprocidad:eq:alpha-inclusion-signos-completa}
\end{equation}
La borne~\eqref{x_reciprocidad:eq:alpha-cota-cola-uniforme} permet de choisir une
troncature qui conserve le signe à chaque extrémité où la fonction complète
ne s'annule pas. Si une extrémité coïncide avec la racine, on utilise l'identité de
la fonction complète et l'appartenance déjà démontrée ; on n'attribue pas ce même
signe à ses troncatures. En effet, le reste exact est
\[
 E_{\mathsf s,\infty}(x)-E_{\mathsf s,M}(x)
 =\lambda(\mathsf s)x^{10}
   \frac{(qx^3)^{M+1}}{1-qx^3}.
\]
Dans le domaine positif considéré, où \(\lambda(\mathsf s)<0\), on obtient à la racine
\[
 E_{\mathsf s,M}(a(\mathsf s))
 =-\lambda(\mathsf s)a(\mathsf s)^{10}
   \frac{(qa(\mathsf s)^3)^{M+1}}{1-qa(\mathsf s)^3}>0
 \qquad(M<\infty).
\]
La certification des signes au sens large dans \eqref{x_reciprocidad:eq:alpha-inclusion-signos-completa} se rapporte donc à \(E_{\mathsf s,\infty}\). Son évaluation utilise la fonction complète ou la troncature accompagnée du reste signé. Cela couvre aussi bien l'infimum appartenant au cylindre que le supremum éventuellement exclu de celui-ci.

\begin{proposition}[Carrés de compatibilité et cylindres communs]
\label{x_reciprocidad:prop:alpha-cuadrados-compatibilidad}
Les restrictions des mots et des publications analytiques commutent :
\[
\begin{array}{ccc}
 \mathcal X_\alpha^{\mathrm{enr}}&
 \xrightarrow{\ \mathsf T_{\alpha,N'}\ }&\mathcal W_{N'}\\[1mm]
 \Big\Vert&&\Big\downarrow\rho_{N',N}\\[1mm]
 \mathcal X_\alpha^{\mathrm{enr}}&
 \xrightarrow{\ \mathsf T_{\alpha,N}\ }&\mathcal W_N
\end{array}
\qquad
\begin{array}{ccc}
 \mathcal X_\alpha^{\mathrm{enr}}&
 \xrightarrow{\ \Gamma_{M'}\ }&\mathfrak J_{M'}\\[1mm]
 \Big\Vert&&\Big\downarrow\tau_{M',M}\\[1mm]
 \mathcal X_\alpha^{\mathrm{enr}}&
 \xrightarrow{\ \Gamma_M\ }&\mathfrak J_M .
\end{array}
\]
La famille \(I_{B,N}\) fournit la connexion effective entre les deux
carrés et satisfait matériellement
\(I_{B,N}\subseteq C_N^A\) pour tout \(N\).
\end{proposition}

\begin{proof}
Le premier carré est la compatibilité de la normalisation entière ; le second est la proposition~\ref{x_reciprocidad:prop:alpha-naturalidad-recurrencia-explicita}. L'inclusion découle de la définition~\eqref{x_reciprocidad:eq:alpha-intervalos-b-completos}, et la non-vacuité et les signes découlent du théorème~\ref{x_reciprocidad:thm:alpha-raiz-completa-explicita}. Aucune étape n'identifie la racine d'une troncature à celle d'une queue arbitraire.
\end{proof}

Nous définissons maintenant, sans laisser implicite le lecteur de préfixes,
\begin{equation}
 \alpha_B(\mathsf s):=
 \operatorname*{arg\,zero}_{x\in I_\alpha}E_{\mathsf s,\infty}(x),
 \qquad
 \operatorname{pref}_N(x):=
 \frac{\lfloor1000^Nx\rfloor}{1000^N},
 \qquad
 \alpha_{B,N}:=\operatorname{pref}_N(\alpha_B).
 \label{x_reciprocidad:eq:alpha-definicion-b-y-prefijos-rev08}
\end{equation}

\begin{theorem}[Égalité projective des deux publications corrélées]
\label{x_reciprocidad:thm:alpha-dos-publicaciones-completas}
Pour tout \(N\), la voie entière et la voie analytique publient le même préfixe :
\begin{equation}
 \boxed{\alpha_{A,N}=\operatorname{pref}_N(\alpha_A)
 =\operatorname{pref}_N(\alpha_B)=\alpha_{B,N}.}
 \label{x_reciprocidad:eq:alpha-dos-vias-cada-nivel}
\end{equation}
À la limite,
\begin{equation}
 \boxed{\alpha_A
 =\operatorname*{arg\,zero}_{x\in I_\alpha}
       E_{\mathsf s,\infty}(x)
 =\alpha_B=\alpha_{\mathrm{HMT}}.}
 \label{x_reciprocidad:eq:alpha-dos-vias-limite-coinductivo}
\end{equation}
\end{theorem}

\begin{proof}
La normalisation entière publie exactement \(a(\mathsf s)\). Le
théorème~\ref{x_reciprocidad:thm:alpha-raiz-completa-explicita} démontre que la carte
analytique représente ce même point comme unique racine simple. Les
cylindres semi-ouverts \(C_N^A\) fixent un préfixe canonique et
\(I_{B,N}\subseteq C_N^A\) contient \(\alpha_B=\alpha_A\). Leurs
préfixes coïncident donc pour chaque \(N\) ; l'emboîtement et les diamètres tendant vers
zéro donnent l'égalité des limites. La reconnaissance
métrologique intervient ensuite.
\end{proof}

La fonction analytique constitue une publication corrélée du même
état, et non une seconde dérivation causalement indépendante de la clôture
entière. Cette précision ne diminue pas l'égalité démontrée : elle identifie son
mécanisme exact, empêche d'attribuer au jet une queue qu'il ne détermine pas et évite la
promotion d'une valeur conventionnelle au rang de générateur.

% END OWNER


% END OWNER

